\begin{abstract}
% Competing entities can often benefit mutually by sharing resources, such as datasets  or access to experimental equipment for research and 
% development. However, uncertainty about an agent's valuation of another's resources can make participants reluctant to share 
% even when mutual benefit is likely, due to fears of imbalanced realized benefits. For example, the value of access to a competitor's 
% dataset or time on experimental equipment is typically revealed only through use and is difficult to assess \emph{ex ante}. In this 
% work, we study the design of exchange algorithms for two competing agents whose valuations are uncertain and revealed incrementally 
% through interaction. We focus on resources whose value depends on the amount shared (e.g. volume of data, hours of equipment use) 
% and is inherently non-retractable (i.e. once data is transmitted or time on equipment is used, it cannot be reclaimed).

% We introduce multi-round \emph{exchange algorithms} that propose sequences of increasing exchanges, leveraging information 
% from prior rounds to improve expected welfare. Our first and contribution is an optimal multi-exchange algorithm under minimal 
% assumptions on valuations, utilities, and welfare functions, obtained by reducing the design problem to an optimal stopping 
% problem over a lattice. The resulting algorithm is also individually rational for reasonable choices of welfare functions.
% %
% We then study a significantly more difficult setting in which agents may exit early. Under this model,
% we define \emph{incentive-compatible} exchange algorithms, to be
% those in which no agent can benefit from leaving prematurely.
% Under added assumptions on the environment, we construct an exchange algorithm which is optimal
% among all incentive-compatible exchange algorithms.


% ---- new

Competing entities can often mutually benefit by sharing resources, such as datasets or access to experimental equipment. 
However, agents are often uncertain about the value they will realize from the others' resources.
For example, the value of access to a competitor's dataset or time on experimental equipment is typically revealed only 
through use and is difficult to assess \emph{ex ante}. This uncertainty can make agents reluctant to share even when mutual 
benefit is likely, due to fears of imbalance in realized benefits; an agent may be concerned that a competitor will 
derive more value from her resources than she would from theirs.

In this work, we study the design of exchange algorithms for two competing agents whose valuations for the other's 
resources are uncertain and only revealed incrementally through interaction. We focus on resources whose value 
depends on the amount shared (e.g. volume of data, hours of equipment use) and is inherently non-retractable 
(i.e. once data is transmitted or time on equipment is used, it cannot be reclaimed). We introduce multi-round 
\emph{exchange algorithms} that propose sequences of increasing exchanges, leveraging information from prior rounds 
to improve expected welfare. Our first contribution is an optimal and individually rational multi-round exchange 
algorithm, obtained by reducing this problem to an optimal stopping problem over the lattice $\NN^2_{\geq 0}$. 
We then study the 
setting in which agents may exit early, defining \emph{incentive-compatible} exchange algorithms to be those
in which no agent can benefit from leaving prematurely.
Under additional mild assumptions on the environment, we construct an exchange algorithm which is optimal
among incentive-compatible exchange algorithms.

\end{abstract}